%% 第26回編集委員会（物理第4巻）：組 電W7b
%% p4-16.tex  電-16　読みもの　電気を届ける・電気を変える（★2〜5。読みもの。3段にしない）
%% 中身は計画書 tools/phys4-plan.md の B「電-16」・A89（系統の周波数と需給、整流の平均値）・A19（交流を直流に変える）。
%%   読みもの（例題は「見本の計算」4題。じぶんでとこう・章末問題はなし）。第3巻 p3-14.tex・第2巻 p2-16.tex と同じ形。
%%   中学・高校の公式を本当に導く所がないので「この章で導く中学・高校の公式」の欄は置かない（計画書 0 の1・K-8）。
%% 公式は導き直さず、持ち主の章を指した（書き上がった原稿の節の番号で確かめた）：
%%   回転の運動エネルギー ½Iω²（力-6 ⑧ ph:力-6:energy。この章では電流 I とぶつけないため、電-12 と同じく J と書いた）、
%%   同期速度 n_s＝120f/p と「発電機の回る速さがそのまま周波数」（電-12 ⑧）、インバーター・PWM・V/f 一定（電-12 ⑨）、
%%   発電所から家までの電圧の段・単相3線式・架空地線（電-11 ①）、送電の損失 R(P/V)²（電-11 ⑥）、三相（電-11 ⑦⑧）、
%%   実効値（電-10 ②）、交流が使われるわけ（電-10 ①）、高調波（電-10 ⑪）、平均の値（5-12⑧）、
%%   過渡現象とサージ（電-9 ⑪・電-13）、百分率インピーダンスと短絡電流（電-13 ⑤）、電磁ノイズ（電-15 ⑫）、
%%   電気の安全（感電・漏電・コードの熱）は電-4 ⑫ を指し、この章の囲みは送電線・電柱・災害のときの安全にしぼった（重ねない）。
%%   揚水発電の位置エネルギー（1 m³・500 m で約 1.4 kWh）と「安く長くためる方法はまだわからない」は熱-16 ⑤・まだわかっていないこと を指した。
%%   設備利用率は熱-16 ⑤。一次遅れ（調速機のモデル）は波-4。
%% ほかの章が電-16 を指している所（grep -n '電-16' tex/p4-*.tex で確かめ、その中身をこの章のその節に置いた）：
%%   ・電-10 ①「東と西で周波数が違うわけは、電-16 で学びます」（頭のコメントは「電-16 ③」）→ ③。
%%   ・電-10 ①「充電器（ACアダプター）で交流を直流に変えています（電-16）」→ ⑤（整流・平滑・DC-DC）。
%%   ・電-11 ⑥ のモデル化の札「整流器やインバーターのように、電流の波の形がくずれる負荷（電-16）」→ ⑤（高調波）。
%%   ・電-11 のこの先の学び「周波数の制御：電-16」「直流送電・半導体のスイッチで電力の流れを変える装置（電-16）」→ ②・⑥。
%%   ・電-12 ⑧「発電機が回る速さが、そのまま交流の周波数になります（電-16）」→ ②。
%%   ・電-12 のこの先の学び「パワーエレクトロニクス（電-16）」→ ⑤。
%%   ・電-9 のこの先の学び「電力系統の守り方は、電-16 で読みます」、電-13 のこの先の学び「サージ…変電所の機器を守る設計（電-16）」→ ⑨。
%%   ・電-15 ⑫「モーターを速く切りかえる回路（電-16 のインバーター）」→ ⑤（電磁ノイズの一言）。
%%   ・電-5 のコメント「配電線の『数千 V』は電-0・電-16 の持ち主に合わせる」→ ① は電-11 ① の 6.6 kV を指した。
%%   ・電-15 の最後のふきだし「最後の電-16 では、電気を作って家まで届けるしくみを、まとめて読んでいく」→ 全体。
%% 前の巻との約束（計画書 A-6）：P20（熱-16 ⑤「つくる量と使う量を合わせるしくみ・東 50 Hz・西 60 Hz は電-16」）→ ②③⑦⑧。
%%   熱-16 のこの先の学び「電力系統のしくみは電-16」→ ①②。
%% 記号（この章。第4部なので面積は A（Y-4））：周波数 f（基準 f_0）、発電と使う電力 P_発・P_使、回るものの運動エネルギー K、
%%   慣性モーメント J、角速度 ω、慣性定数 H〔s〕（第2巻のエンタルピー H、電-7 の磁場の強さ H とは別。記号カード）、
%%   つながっている発電機の定格の出力の合計 P_n、面積 A、効率 η（電-15 の波のインピーダンスとは別の章）、電圧 V・電流 I（直流・実効値）、
%%   瞬時の電圧 v（電-10 と同じ）。デューティ比は記号を作らず「オンの割合」とことばで書いた。
%%   第27回（修4-6・G-5）：調速機の調定率は、抵抗 R とぶつけないため記号を作らず「調定率」とことばで書き、式は数 0.05 のまま書いた
%%   （プログラムの変数名は droop）。負荷が周波数で減る係数 D（巻末の「分野で記号がぶつかるもの」の D の行）、発電機の電圧 E_1・E_2、
%%   位相角 δ（もっと知りたい・専門書）、日光の強さ G（例題3）、運動エネルギー K は、この章だけの意味。
%%   図3・図4のラベルの白地をやめた（G-8）。記録は tools/FIX-phys4-review-6.md。
%% 数値はすべて Python で計算した：checks/物理第4巻/電-16/calc_p4_16.py（出力 calc_p4_16_out.txt）、
%%   ② のプログラムは checks/物理第4巻/電-16/sim_book_p4_16.py（出力 sim_book_p4_16_out.txt。本文の実行結果と同じ）。
%% 出典（手引き BRIEF-phys4 2章の13・16。11月5日に書くときに確かめた。Wikipedia は出典にしていない。
%%   電力広域的運営推進機関・東京電力・経済産業省・資源エネルギー庁・電気学会のサイトは作業環境から届かない（occto は接続できず、meti は 403）。
%%   届かない資料は、WebSearch の抜粋で2つ以上を照らし合わせた（「照合」）。
%%   ・周波数を標準周波数に維持するよう努める：電気事業法 第26条第1項、電気事業法施行規則 第38条第2項（「経済産業省令で定める周波数の値は、
%%     その者が供給する電気の標準周波数に等しい値とする」）。e-Gov 法令 API（laws.e-gov.go.jp/api/1/lawdata/339AC0000000170・407M50000400077）を curl で読んだ。
%%   ・周波数の調整の目標（東北・東京 50±0.2 Hz、中部〜九州 60±0.2 Hz、北海道 50±0.3 Hz、沖縄 60±0.3 Hz）：電力広域的運営推進機関「電気の質に関する報告書」
%%     （各年度。2024年度版は 2026年2月）の抜粋、東京電力パワーグリッド「周波数調整・需給運用ルール」（2022年4月改定。50.0±0.2 Hz）の抜粋、
%%     電気学会誌 128巻7号 431 ページ（2008年。J-STAGE の PDF を curl で読んだ。50 Hz・0.2 Hz、1993〜2002年に 99.999 % の数を確かめた）で照合。
%%     紙面は「東京などでは 50 Hz の上下 0.2 Hz、北海道では上下 0.3 Hz」と書いた。
%%   ・東 50 Hz・西 60 Hz の由来（明治の半ば、1890年代に、東京の電灯会社がドイツ製の 50 Hz、大阪の電灯会社がアメリカ製の 60 Hz の発電機を入れた）：
%%     パナソニック「周波数」の解説、ベネッセ教育情報サイト（2021年3月）、hatsudenki.jp の解説で照合。入れた年は資料で 1895〜1896年と食い違うので、
%%     紙面は「1890年代」とだけ書いた。境目（富士川・糸魚川）は電-10 ① と同じ。
%%   ・東西の連系：2021年3月31日に岐阜県と長野県を結ぶ周波数変換設備（90万 kW）が運転を始め、東西の連系の容量は 120万 kW から 210万 kW へ：
%%     関西電力送配電の報道資料（2021年4月1日）、東芝エネルギーシステムズの発表（2021年4月1日）、電気新聞（2021年）、東京電力パワーグリッドの発信で照合。
%%     300万 kW への増強：2021年の時点では 2027年度の運用開始をめざしていたが、2026年5月の広域系統整備委員会の資料（抜粋）で「運開期が遅延する見込み」。
%%     紙面は「300万 kW に増やす工事が進められている（完成の時期は見直されている）」と書いた。
%%   ・北海道と本州の直流の連系：既設 60万 kW に、2019年3月28日に運転を始めた新北本連系設備 30万 kW（直流 250 kV・1200 A、自励式、津軽海峡の下の
%%     鉄道のトンネルを通るケーブル）が加わり 90万 kW：電気学会誌 139巻8号 556 ページ（2019年、北海道電力 三浦）、電気学会論文誌 D 141巻6号
%%     NL6_6（2021年）。どちらも J-STAGE の PDF を curl で読んだ。「日本で初めて電力系統が全域停電するブラックアウト」も前者の文。
%%   ・2018年9月6日の北海道の全域停電：3時7分に地震（M6.7、最大震度7）、震源の近くの石炭火力の2基がタービンの振動を検知して止まり（約116万 kW）、
%%     送電線4回線の事故で水力も止まった。周波数が急に下がり、北本連系の自動周波数制御（49.62 Hz で動作）と周波数低下リレーで計約130万 kW の負荷を
%%     自動で切りはなした。残った1基も止まり、3時25分に全域停電（最大約295万戸）：電力広域的運営推進機関の検証委員会の資料（中間報告 2018年10月、
%%     最終報告 2018年12月）の抜粋、北海道新聞の記事、生活協同組合研究 2019年3月（J-STAGE。M6.7）の抜粋で照合。
%%     発電所の名前に常用漢字にない字（苫）があるので、紙面は「震源の近くにあった大きな石炭火力発電所」と書いた。
%%   ・大停電の大きさの分布：Carreras, Dobson らが北米の記録から「べき乗則に近い」ことを示し、自己組織化臨界の考えを出した。これに対して、
%%     「とても小さいかとても大きいかの2つに分かれる（2つの山）」という研究もあり、決着していない：Kornbluth ら arXiv:2008.01141 の要旨（curl で読んだ）、
%%     Carreras ら IEEE Trans. Circuits Syst. I 51, 1733 (2004)（PSERC・UVM の写し）と Dobson ら Chaos 12, 985 (2002) の抜粋で照合。
%%   ・再生可能エネルギーの割合：2024年度の発電電力量 9911億 kWh、再エネ（水力をふくむ）23.1 %（資源エネルギー庁「令和6年度エネルギー需給実績（確報）」
%%     2026年4月14日。経済産業省の報道発表の抜粋とコージェネ財団の要約で照合）。2023年度の太陽光 9.8 %（同 2023年度確報、eneres・日本経済新聞で照合）。
%%   ・太陽光の出力制御：2018年10月に九州（離島を除く）で国内初（SOLAR JOURNAL、sustainablejapan、電力広域的運営推進機関の系統 WG 資料の抜粋で照合）。
%%   ・太陽電池の公称最大出力は 1000 W/m²・AM1.5・25 ℃の条件、モジュールの変換効率は 20 % 前後：パナソニック・京セラの製品の説明の抜粋で照合。
%%   ・慣性定数：火力の 2H が約8 s（H 約4 s）、複合発電 2H 約11 s：日本エネルギー学会誌 44巻4号 171（J-STAGE）の抜粋。日本エネルギー経済研究所の資料、
%%     電力広域的運営推進機関の資料（2020年10月）で「インバーターの電源が増えると慣性が小さくなり、周波数の変わる速さが大きくなる」ことを照合。
%%     見本の計算は「H＝5 s とします」と書いた。
%%   ・揚水発電：全国 約2700万 kW（2022年の時点。資源エネルギー庁「エネこれ」の抜粋、経済産業省の勉強会資料 2750万 kW（2020年度）で照合）、
%%     総合効率 約70 %（JICA の調査報告、資源エネルギー庁の「くみ上げで約3割のロス」で照合）。
%%   ・自励式と他励式の変換器、ブラックスタート：上の電気学会誌・論文誌の2つ。
%%   ・切れた電線・送電線の近くの安全：近畿産業保安監督部（経済産業省の地方の部局）の案内、東北電力・東京電力パワーグリッドの案内、
%%     電気事業連合会の案内（壊れた太陽光パネルも日が当たると発電する）の抜粋で照合。避難のときにブレーカーを切ること：関西電力送配電の案内の抜粋
%%     （内閣府などの防災の案内でも広く言われている）。
%%   ・〔要確認〕なし（見本の計算の H＝5 s・送電線の抵抗 2 Ω・調速機の数は仮の値として本文で「〜とします」と書いた）。
\makeatletter
% 箱  \sxqbox{名前}{(位置)}{はば(mm)}{字}{色}
\newcommand\sxqbox[5]{\node[draw=#5!75!black, fill=#5!10, line width=0.8pt, rounded corners=1.4mm, minimum width=#3 mm, minimum height=10mm,
  text width=#3 mm-4mm, align=center, font=\sffamily\footnotesize, inner sep=2pt] (#1) at #2 {#4};}
% 小さな家  \sxqhouse{(左下)}
\newcommand\sxqhouse[1]{\begin{scope}[shift={#1}]
  \filldraw[fill=bkOrange!15, draw=bkInk, line width=0.6pt] (0,0) rectangle (7,5);
  \filldraw[fill=bkRed!35, draw=bkInk, line width=0.6pt] (-0.8,5) -- (3.5,8.5) -- (7.8,5) -- cycle;
  \fill[bkYellow!70] (2.6,1.2) rectangle (4.4,3.2);\end{scope}}
% 発電所（煙突つき）  \sxqplant{(左下)}
\newcommand\sxqplant[1]{\begin{scope}[shift={#1}]
  \filldraw[fill=bkGray!12, draw=bkInk, line width=0.6pt] (0,0) rectangle (11,7);
  \filldraw[fill=bkGray!12, draw=bkInk, line width=0.6pt] (7.5,7) rectangle (9.5,16);
  \fill[bkGray!40] (8.6,17.4) circle (1.0) (9.8,19.2) circle (1.4);\end{scope}}
% 鉄塔  \sxqtower{(根もと)}
\newcommand\sxqtower[1]{\begin{scope}[shift={#1}]
  \draw[bkInk, line width=0.7pt] (-3,0) -- (0,18) -- (3,0) (-2.2,5) -- (2.2,5) (-1.4,10) -- (1.4,10) (-4,14) -- (4,14) (-3,17) -- (3,17)
    (-2.2,5) -- (1.4,10) (2.2,5) -- (-1.4,10);\end{scope}}
\makeatother
\tikzset{sxq lab/.style={font=\sffamily\footnotesize, inner sep=1pt, text=bkInk},
  sxq arr/.style={-{Stealth[length=2.4mm, width=2mm]}, line width=1pt, bkInk!80},
  sxq info/.style={-{Stealth[length=2mm, width=1.8mm]}, line width=0.7pt, bkBlue!80!black, densely dashed},
  sxq dim/.style={bkGray!80!black, line width=0.5pt, {Stealth[length=1.8mm]}-{Stealth[length=1.8mm]}}}

\phchapter[読みもの]{16}{電気を届ける・電気を変える}{停電しないように、電気の量をどうやって合わせている？}{\bklevelrange{2}{5}}{電気を届ける・電気を変える}
\bkmichishirube{力-6,熱-16,電-4,電-9,電-10,電-11,電-12,電-13,電-15}{原-9,原-10,原-15}

\begin{goal}
\begin{itemize}[label={\color{bkTheme}●}, leftmargin=1.6\zw, itemsep=0pt, topsep=1pt]
\item \bklevel{2}\enspace 発電所から家までの電気の道すじと、東と西で周波数が違うわけ・東西をつなぐしくみを説明できる。
\item \bklevel{3}\enspace 整流・インバーター・直流送電を、スイッチの切りかえと「平均の値」で説明し、整流した電圧の平均を計算できる。
\item \bklevel{4}\enspace つくる電気と使う電気がずれると周波数が変わるわけを、回る発電機の運動エネルギーで説明し、周波数の下がる速さを見積もれる。
\item \bklevel{5}\enspace 調速機の働きを入れた周波数の変わり方を、プログラムで計算して読める。
\end{itemize}
\end{goal}

\begin{hatena}
\noindent\begin{minipage}[c]{0.55\linewidth}\raggedright
夕方、たくさんの家で一斉に明かりやエアコンをつけても、電気はふつう止まりません。でも、電気は大きな池にためておいて、使う分だけくみ出しているのではありません。発電所でつくったその瞬間に、どこかで使われています。\par\vspace{3pt}
使う量は、何千万もの家や工場が、それぞれ好きなときに勝手に変えています。発電所は、それをどうやって知り、つくる量をぴったり合わせているのでしょう。
\end{minipage}\hfill
\begin{minipage}[c]{0.42\linewidth}\centering
\begin{tikzpicture}[x=0.9mm, y=0.9mm]
  \fill[bkGray!20] (-2,-2) rectangle (56,0);
  \draw[bkInk, line width=0.6pt] (-2,0) -- (56,0);
  \sxqplant{(0,0)}
  \sxqtower{(23,0)}
  \draw[bkInk!70, line width=0.5pt] (11,6) .. controls (16,9) .. (19,14) (27,14) .. controls (33,10) .. (38,6);
  \sxqhouse{(38,0)}
  \sxqhouse{(47,0)}
  % 天びん
  \draw[bkInk, line width=0.8pt] (27,24) -- (27,30);
  \draw[bkInk, line width=0.9pt] (15,31) -- (39,29);
  \fill[bkInk] (27,30) circle (0.8);
  \filldraw[fill=bkBlue!25, draw=bkBlue!70!black, line width=0.5pt] (11,28) rectangle (19,31);
  \filldraw[fill=bkOrange!30, draw=bkOrange!70!black, line width=0.5pt] (35,26) rectangle (43,29);
  \node[sxq lab, anchor=south] at (15,31.5) {つくる};
  \node[sxq lab, anchor=south] at (39,29.5) {使う};
  \node[font=\sffamily\bfseries\Large, text=phHatena!80!black] at (52,22) {?};
\end{tikzpicture}
\end{minipage}
\end{hatena}

この章は読みものです。じぶんでとこう・章末問題はありません（見本の計算の例題はあります）。これまでの章で学んだ「回る物のエネルギー」「発電機と変圧器」「三相交流と送電」「交流の実効値と平均」「過渡現象」などの道具が、国じゅうに広がる電気の網の中で、どう使われているかを見ていきます。先に結論を言うと、電気の網の中では、\textbf{周波数}（\qty{50}{Hz}・\qty{60}{Hz}）が「つくる量と使う量のつり合い」を知らせる目もりになっています。この章で使う、これまでの章の道具を、先に見ておきましょう（図\ref{ph:電-16:fig-michi}）。
\begin{bkfigure}
\begin{tikzpicture}[x=1mm, y=1mm]
  \tikzset{src/.style={bk node, text width=22mm, inner xsep=1mm, align=center, minimum height=17mm, font=\sffamily\footnotesize}}
  \node[src] (a) at (0,24) {回る物の\\エネルギー\\（力-6）};
  \node[src] (b) at (27.2,24) {発電機・\\インバーター\\（電-12）};
  \node[src] (c) at (54.4,24) {三相交流と\\送電の損失\\（電-11）};
  \node[src] (d) at (81.6,24) {実効値・\\高調波\\（電-10）};
  \node[src] (e) at (108.8,24) {過渡現象・\\サージ\\（電-9・電-13）};
  \node[bk node, draw=bkTheme, line width=1pt, text width=96mm, align=left, minimum height=11mm, font=\sffamily\small] (f) at (54.4,-2) {電気を届け、形を変え、つり合いを保つ（電-16）};
  \foreach \n in {a,b,c,d,e} {\draw[bk link] (\n.south) -- (\n.south |- f.north);}
\end{tikzpicture}
\bkfigcaption{これまでの章の道具が、電気を届けるしくみの中で使われている。}\label{ph:電-16:fig-michi}
\end{bkfigure}

\begin{yattemiyou}[電気製品のラベルと「でんき予報」を見よう]
\begin{enumerate}[label={\bkcircled[bkGreen]{\arabic*}}, leftmargin=2.4\zw, itemsep=1pt, topsep=2pt]
\item 電子レンジ・ドライヤー・充電器などの、本体や底に貼ってあるラベル（定格の表示）を見る。「\qty{50}{Hz}」「\qty{60}{Hz}」「50/60 Hz」のどれが書いてあるかを調べる。ラベルを見るときは、電源プラグをぬいてから、大人といっしょに見る。中を開けたりしない。
\item 充電器（ACアダプター）のラベルの「入力 AC \qty{100}{V}」「出力 DC \qty{5}{V}」のような表示を書きとめる。AC は交流、DC は直流のことである。
\item 住んでいる地域の電力会社のウェブページで、「でんき予報」（今日の電気の使用率のグラフ）を見て、1日のうち、使う電気がいちばん多い時刻を調べよう。
\end{enumerate}
\end{yattemiyou}

\noindent 多くの充電器は「AC 100〜\qty{240}{V}、50/60 Hz」と書いてあり、東でも西でも、外国でも使えます。中で、いったん交流を直流に変えてしまうからです（⑤）。一方、古い電気製品には、\qty{50}{Hz}用と\qty{60}{Hz}用が分かれているものがあります。なぜ日本には2つの周波数があるのかは③で、使う電気が多い時刻に、電力会社がどうやって電気を足しているのかは②で見ます。

\begin{omoidasou}
\bkrecall{力-6}{回る物の運動エネルギーは $\tfrac12I\omega^2$（⑧）。この章では、電流 $I$ とぶつけないように、慣性モーメントを電-12 と同じ $J$ と書く。}
\bkrecall{電-12}{同期発電機の回る速さは $n_{\mathrm s}=120f/p$（⑧）。2極で\qty{50}{Hz}なら毎分3000回転。インバーターはスイッチの切りかえで直流から交流をつくる（⑨）。}
\bkrecall{電-11}{同じ電力 $P$ を電圧 $V$ で送ると、電線の損失は $R(P/V)^2$（⑥）。発電所から家までの電圧の段（①）。}
\bkrecall{5-12}{平均の値は、グラフの下の面積を全体のはばでわったもの $\dfrac{1}{b-a}\displaystyle\int_a^bf(t)\,\dd t$（⑧）。}
\end{omoidasou}

%% ================================================================
\Needspace{16\baselineskip}
\bksection{doguu}{物理の道具}

\bksubsection{1}{電気が届くまで：国じゅうが1つの機械}
\leavevmode\label{ph:電-16:sec1}発電所でつくった電気は、変電所で電圧を何回も変えながら、家や工場まで届きます。送電線は\qty{500}{kV}・\qty{275}{kV}・\qty{154}{kV}・\qty{66}{kV}などの三相交流、町の配電線は\qty{6.6}{kV}の三相交流、家の中は単相の\qty{100}{V}・\qty{200}{V}です（電-11 ①の図）。高い電圧で送るのは、同じ電力なら電流が小さくなり、電線の損失 $R(P/V)^2$ が小さくなるからでした（電-11 ⑥）。

大切なのは、この道すじが1本の線ではなく、\textbf{網の目}になっていることです。たくさんの発電所と、たくさんの家や工場が、送電線の網でつながっています。1つの送電線や発電所が故障しても、ほかの道を通って電気が届くようにするためです。このように、発電所・送電線・変電所・配電線と、電気を使う人までをつないだ全体を、\textbf{電力系統}（でんりょくけいとう）といいます（図\ref{ph:電-16:fig-keitou}）。\bkidx{でんりょくけいとう}{電力系統}
\begin{bkfigure}
\begin{tikzpicture}[x=1mm, y=1mm]
  % 修4-6（E6-22）：①の本文のページに入るよう、たてを少しつめた
  \sxqbox{g1}{(11,26)}{22}{火力}{bkRed}
  \sxqbox{g2}{(11,15)}{22}{水力・揚水}{bkBlue}
  \sxqbox{g3}{(11,4)}{22}{太陽光・風力}{bkGreen}
  \sxqbox{s1}{(46,23)}{20}{変電所}{bkGray}
  \sxqbox{s2}{(46,7)}{20}{変電所}{bkGray}
  \sxqbox{s3}{(80,15)}{20}{配電用変電所}{bkGray}
  \sxqbox{u1}{(113,25)}{20}{工場・ビル}{bkOrange}
  \sxqbox{u2}{(113,5)}{20}{家}{bkOrange}
  \draw[bkInk, line width=1.2pt] (g1) -- (s1) (g2) -- (s1) (g2) -- (s2) (g3) -- (s2) (s1) -- (s2) (s1) -- (s3) (s2) -- (s3) (s3) -- (u1) (s3) -- (u2);
  \sxqbox{c}{(62,-7.5)}{76}{中央給電指令所（周波数と電力の流れを見張る）}{bkPurple}
  \draw[sxq info] (c.west) -| (g3.south);
  \draw[sxq info] (c.east) -| (u2.south);
  \node[sxq lab, anchor=south] at (28,31.5) {送電線（三相）};
  \node[sxq lab, anchor=south] at (98,22) {配電線};
\end{tikzpicture}
\bkfigcaption{電力系統（形を見るための図）。太い線が電気の通る網、破線は知らせ（周波数・電力の量・指令）の流れ。どこかの線が切れても、ほかの道で届くように網にしてある。}\label{ph:電-16:fig-keitou}
\end{bkfigure}

網でつながっている、ということは、もう1つ大事な意味があります。同じ網につながったすべての交流の発電機は、\textbf{同じ周波数で、足並みをそろえて回っている}のです。電-12 ⑧で見たとおり、同期発電機は、回る速さと交流の周波数が決まった関係（$n_{\mathrm s}=120f/p$）でしばられています。網につながった発電機は、たがいに電気でつながって引っぱり合い、1台だけ速く回ったり遅く回ったりできません。東日本の何百台もの発電機は、1つの大きな機械の部品のように、そろって\qty{50}{Hz}で回っています。

\begin{motto}[5]
\phboxhead{足並みをそろえる力}
2台の発電機を、リアクタンス $X$ の送電線（抵抗は無視する）でつなぎます。2台の電圧の大きさを $E_1$・$E_2$、位相のずれ（回転子の角のずれ）を $\delta$ とすると、電-10 のフェーザで計算して、送電線を流れる電力は
\[P=\frac{E_1E_2}{X}\sin\delta\]
です。発電機1が少し先に進むと、$\delta$ が大きくなり、発電機1が送り出す電力が増えます。電力を多く出すほど回転にはブレーキがかかり、受けとる側の発電機2は押されて速くなるので、ずれはもとにもどります。網の上の発電機は、この「ずれると引きもどす」力で、足並みをそろえているのです。ただし、$\sin\delta$ は $\delta=90^\circ$ でいちばん大きくなるので、それより多くの電力は送れません。大きな故障で $\delta$ が大きくずれすぎると、引きもどせなくなり、発電機は足並みからはずれてしまいます（脱調）。
\end{motto}

\begin{phmijika}[中央給電指令所]
電力会社には、系統全体を見張る\textbf{中央給電指令所}があります。大きな画面に、いまの周波数、それぞれの発電所の出力、送電線を流れる電力が表示され、昼も夜も、人とコンピュータが見張っています。前の日には、天気予報や曜日から、次の日の使う量を予想して、どの発電所を何時に動かすかを決めておきます。
\end{phmijika}

%% ================================================================
\bksubsection{2}{周波数で見張る：はてな？の答え}
\leavevmode\label{ph:電-16:sec2}電気は、ためておかずに、つくった瞬間に使われます。それなら、使う量が急に増えたとき、足りない分はどこから来るのでしょう。

答えは、\textbf{回っている発電機そのもの}です。火力・原子力・水力の発電所のタービンと発電機は、何十トン、何百トンもある重い回転体で、毎分3000回転（\qty{50}{Hz}・2極）などの速さで回っています。回る物は、運動エネルギー $K=\tfrac12J\omega^2$ をもっています（力-6 ⑧）。使う電力 $P_{\text{使}}$ が、タービンが発電機に入れる電力（つくる電力）$P_{\text{発}}$ より大きくなると、足りない分は、この回る運動エネルギーから出ていきます。
\pheq{\frac{\dd K}{\dd t}=P_{\text{発}}-P_{\text{使}}}{J/s ＝ W}
\label{ph:電-16:swing}運動エネルギーが減れば、回転は遅くなり、周波数が下がります。反対に、使う量が減ると、回転が速くなり、周波数が上がります。

\begin{phyomi}[$\dfrac{\dd K}{\dd t}=P_{\text{発}}-P_{\text{使}}$]
「つくる量と使う量の差だけ、回る発電機のエネルギーが増えたり減ったりする」。これは、はずみ車（力-6 ④）と同じエネルギーの保存の式です。周波数は、はずみ車の回る速さ、つまり系統にたまっているエネルギーの目もりです。周波数が下がっていれば「つくる量が足りない」、上がっていれば「つくりすぎ」とわかります。
\end{phyomi}

\phboxhead{周波数の下がる速さ}
発電機の回る速さ $\omega$ は周波数 $f$ に比例するので、$K=\tfrac12J\omega^2$ は $f^2$ に比例します。基準の周波数 $f_0$ のときの運動エネルギーを $K_0$ とすると、$K=K_0(f/f_0)^2$ です。これを時間で微分すると（合成関数の微分）、$f$ が $f_0$ に近い間は
\[\frac{\dd K}{\dd t}=\frac{2K_0f}{f_0^{\,2}}\,\frac{\dd f}{\dd t}\approx\frac{2K_0}{f_0}\,\frac{\dd f}{\dd t}\]
です。電力の工学では、回る運動エネルギー $K_0$ を、つながっている発電機の定格の出力の合計 $P_{\mathrm n}$ が何秒分か、で表します。これを\textbf{慣性定数} $H$〔s〕といい、$K_0=HP_{\mathrm n}$ です。\bkidx{かんせいていすう}{慣性定数}火力発電所の発電機では、$H$ は数秒です。
\begin{phflow}[周波数の下がる速さ]
\phflowrow{$\dfrac{\dd K}{\dd t}=P_{\text{発}}-P_{\text{使}}$}{$K=K_0(f/f_0)^2$\\$K_0=HP_{\mathrm n}$}{$\dfrac{\dd f}{\dd t}=\dfrac{f_0}{2H}\cdot\dfrac{P_{\text{発}}-P_{\text{使}}}{P_{\mathrm n}}$}
\end{phflow}
右の式は、「足りない電力が全体の何割か」に $f_0/(2H)$ をかけると、1秒に周波数が何 Hz 変わるかがわかる、という式です。\qty{50}{Hz}の系統で、$H=\qty{5}{s}$ のとき、全体の1\,\% の発電所が急に止まると、周波数は1秒に\qty{0.05}{Hz}の速さで下がり始めます（\bkreidairef{1}）。
\begin{phbaisuru}{慣性定数 $H$}
$H$ を2倍にすると、同じ電源が止まっても、周波数の下がる速さは半分になります。重いはずみ車ほど、同じブレーキでもゆっくりしか遅くならないのと同じです。反対に、⑦で見る太陽光や風力の発電のように、回る重い発電機を持たない電源が増えると、系統の $H$ は小さくなり、周波数は速く変わるようになります。
\end{phbaisuru}

\begin{senmonsho}[電力系統工学の本の書き方]
\begin{bkhonbun}
同期発電機の動揺方程式は $\dfrac{2H}{\omega_{\mathrm s}}\dfrac{\dd^2\delta}{\dd t^2}=P_{\mathrm m}-P_{\mathrm e}$〔p.u.〕で表される．ここで $H$ は単位慣性定数〔s〕，$\delta$ は内部誘導起電力の位相角，$P_{\mathrm m}$ は機械入力，$P_{\mathrm e}$ は電気出力である．
\end{bkhonbun}
\bkyaku{「p.u.」は単位法（電-13 ⑤）で、電力を定格の何倍かで書く。$\delta$ は発電機の回転子の位置（角）を、系統の基準からどれだけ進んでいるかで表した量で、$\dd\delta/\dd t$ が回る速さのずれ（周波数のずれ）にあたる。この章の $\dd K/\dd t=P_{\text{発}}-P_{\text{使}}$ を、1台の発電機について角で書いたもの。}
\end{senmonsho}

\phboxhead{3つの段で、つくる量を合わせる}
回る運動エネルギーは、すぐに足りない分を出してくれますが、出し続ければ周波数は下がり続けます。そこで、電力系統は、時間の長さの違う3つの段で、つくる量を使う量に合わせ直します（図\ref{ph:電-16:fig-dan}）。
\begin{enumerate}[label={\sffamily\bfseries 段\arabic*}, leftmargin=2.4\zw, itemsep=2pt, topsep=2pt]
\item \textbf{慣性}（その瞬間）：回る発電機の運動エネルギーが、何もしなくても足りない分を出す。周波数が下がり始める。
\item \textbf{調速機}（数秒〜数十秒）：発電所のタービンには、回る速さをはかって、遅くなると蒸気や水を多く入れる\textbf{調速機}（ガバナー）がついている。\bkidx{ちょうそくき}{調速機}それぞれの発電所が、自分の回る速さ（＝系統の周波数）だけを見て、自動で出力を増やす。周波数の下がりは止まるが、基準より少し低い所で落ち着く。
\item \textbf{周波数をもとにもどす制御}（数分）：中央給電指令所のコンピュータが、周波数のずれと、ほかの地域との電力のやりとりをはかり、指定した発電所に「出力を増やせ」と指令を送る（負荷周波数制御、LFC）。周波数は基準の値にもどる。さらに長い時間では、燃料の安い発電所から順に出力を割りふり直す（経済負荷配分）。
\end{enumerate}
\begin{bkfigure}
\begin{tikzpicture}
\begin{axis}[bkplot, axis x line*=bottom, axis y line*=left, width=112mm, height=50mm, xmin=0, xmax=60, ymin=-0.21, ymax=0.03, xtick={0,10,20,30,40,50,60}, ytick={-0.2,-0.15,-0.1,-0.05,0}, yticklabels={$-0.20$,$-0.15$,$-0.10$,$-0.05$,0}, xlabel={時刻（s）}, ylabel={周波数のずれ（Hz）}]
  \draw[bkGray, line width=0.5pt] (axis cs:0,0) -- (axis cs:60,0);
  \draw[bkRed!80!black, line width=1.0pt, densely dotted] (axis cs:0,0) -- (axis cs:4,-0.2);
  \draw[bkBlue!80!black, line width=1.5pt] (axis cs:0,0) -- (axis cs:1,-0.0456) -- (axis cs:2,-0.0728) -- (axis cs:3,-0.0761) -- (axis cs:4,-0.0597) -- (axis cs:5,-0.0341) -- (axis cs:6,-0.0107) -- (axis cs:7,0.0026) -- (axis cs:8,0.0033) -- (axis cs:9,-0.0060) -- (axis cs:10,-0.0196) -- (axis cs:11,-0.0316) -- (axis cs:12,-0.0380) -- (axis cs:13,-0.0378) -- (axis cs:14,-0.0326) -- (axis cs:15,-0.0254) -- (axis cs:16,-0.0193) -- (axis cs:17,-0.0162) -- (axis cs:18,-0.0166) -- (axis cs:19,-0.0195) -- (axis cs:20,-0.0233) -- (axis cs:21,-0.0264) -- (axis cs:22,-0.0279) -- (axis cs:23,-0.0275) -- (axis cs:24,-0.0259) -- (axis cs:25,-0.0239) -- (axis cs:26,-0.0224) -- (axis cs:27,-0.0217) -- (axis cs:28,-0.0219) -- (axis cs:29,-0.0228) -- (axis cs:30,-0.0238) -- (axis cs:32,-0.0249) -- (axis cs:34,-0.0243) -- (axis cs:36,-0.0234) -- (axis cs:38,-0.0233) -- (axis cs:40,-0.0239) -- (axis cs:44,-0.0239) -- (axis cs:48,-0.0237) -- (axis cs:52,-0.0239) -- (axis cs:56,-0.0238) -- (axis cs:60,-0.0238);
  \draw[bkGreen!50!black, line width=1.0pt, dashed] (axis cs:30,-0.0238) .. controls (axis cs:42,-0.022) and (axis cs:48,-0.004) .. (axis cs:60,-0.001);
  % 修4-6（G-8）：ラベルは白地をやめ、格子の線とグラフの線のない升目の中に置いた
  \node[sxq lab] at (axis cs:6.5,-0.175) {段1};
  \node[sxq lab] at (axis cs:15,-0.075) {段2};
  \node[sxq lab] at (axis cs:45,0.015) {段3};
\end{axis}
\end{tikzpicture}
\bkfigcaption{1\,\% の発電所が止まったときの周波数（$H=\qty{5}{s}$ の計算）。点線（段1）は慣性だけで、4秒で $-0.2$\,Hz まで下がる。実線（段2）は調速機ありで、下がりが止まる（もっと知りたい「調速機の働きを、プログラムで見る」の計算）。破線（段3）は指令所の制御で基準の値へもどるようすを、形だけかいたもの（実際には数分かかる）。}\label{ph:電-16:fig-dan}
\end{bkfigure}
\noindent 日本では、法令（電気事業法施行規則 第38条）で、周波数は「標準周波数に等しい値」に維持するよう努めることになっています。実際の目標は、東京などの地域で\qty{50}{Hz}の上下\qty{0.2}{Hz}、西日本で\qty{60}{Hz}の上下\qty{0.2}{Hz}、系統の小さい北海道では\qty{50}{Hz}の上下\qty{0.3}{Hz}です。\qty{0.2}{Hz}は、\qty{50}{Hz}の0.4\,\% です。日本の系統は、ふだんは、ほとんどいつも、この目標の中に収まっています。

\begin{motto}[5]
\phboxhead{調速機の働きを、プログラムで見る}\label{ph:電-16:prog}%
全体の1\,\% の発電所が止まったときの周波数のずれを、表計算やプログラムで計算してみます。周波数のずれ（基準の何倍か）を $x$、調速機が増やした出力（全体の何倍か）を $y$ とします。上の式に、使う電力が周波数とともに少し減る分（モーターが遅く回る。$D=1$ とします）を入れると
\[2H\frac{\dd x}{\dd t}=y-0.01-Dx,\qquad T_{\mathrm g}\frac{\dd y}{\dd t}=-\frac{x}{0.05}-y\]
です。2つめの式が調速機で、「周波数が下がると、そのずれの $1/0.05=20$ 倍まで、時間 $T_{\mathrm g}$ ほどかけて出力を増やす」という一次おくれ（波-4）のモデルです。0.05 は\textbf{調定率}といい、「周波数が5\,\% 下がると出力を100\,\% 増やす」割合です。$H=\qty{5}{s}$、$T_{\mathrm g}=\qty{5}{s}$ とします。電-9 ⑩と同じく、短い時間 $\Delta t$ ごとに値を先に進めます。
\begin{phcode}[Python]
f0 = 50.0
H, D, Tg = 5.0, 1.0, 5.0
droop = 0.05    \# 調定率
dPL = 0.01      \# 止まった発電所（全体の 1 \%）
dt = 0.1
df, dPg = 0.0, 0.0   \# x と y
for n in range(301):
    if n <= 40 and n \% 10 == 0 or n \% 100 == 0:
        print(round(n*dt, 1), round(df*f0, 3))
    ddf = (dPg - dPL - D*df) / (2*H)
    ddPg = (-df/droop - dPg) / Tg
    df, dPg = df + ddf*dt, dPg + ddPg*dt
\end{phcode}
\begin{phcodeout}[実行結果（時刻 s と、周波数のずれ Hz）]
0.0 0.0
1.0 -0.046
2.0 -0.073
3.0 -0.076
4.0 -0.06
10.0 -0.02
20.0 -0.023
30.0 -0.024
\end{phcodeout}
はじめの1秒は、\bkreidairef{1}の $\qty{0.05}{Hz/s}$ に近い速さで下がります。調速機が出力を増やすと、約3秒で下がりが止まり（いちばん下がって約 $-\qty{0.08}{Hz}$）、少しゆれてから $-\qty{0.024}{Hz}$ ほどで落ち着きます（図\ref{ph:電-16:fig-dan}）。この残ったずれを0にもどすのが、段3の制御です。
\end{motto}


\phboxhead{それでも足りないとき：自動で一部を切りはなす}
大きな発電所がいくつも同時に止まると、調速機が間に合わず、周波数が下がり続けることがあります。周波数が低いまま回り続けると、タービンの羽根がゆれて傷んだりするので、発電所は自分を守るために系統から切りはなされます。すると、さらにつくる量が減り、周波数がもっと下がる、という悪い循環が起きます。そこで、周波数が決まった値より下がったら、\textbf{一部の地域への送電を自動で切りはなして}、使う量を減らし、残りの系統を守るしくみ（周波数低下リレー）が備えてあります。一部の地域の停電とひきかえに、全体の停電を防ぐのです。

\phboxhead{はてな？の答え}
\bkhatenaref の答え：電力会社は、使う量を1軒ずつはかって合わせているのではありません。\textbf{周波数を見ています}。使う量がつくる量より増えると、回る発電機の運動エネルギーが使われ、系統じゅうの発電機がそろって少し遅くなり、周波数が下がります。それぞれの発電所の調速機が、その周波数の下がりを見て、数秒のうちに自動で出力を増やし、さらに中央給電指令所が、数分かけて周波数を基準の値にもどすように発電所に指令を出します。前の日には、天気や曜日から使う量を予想して、発電所の動かし方を決めておきます。それでも足りないほど大きな発電所が止まったときは、一部の地域への送電を自動で切りはなして、全体の停電を防ぎます。

%% ================================================================
\bksubsection{3}{50 Hz と 60 Hz：東と西をつなぐ}
\leavevmode\label{ph:電-16:sec3}日本では、東日本（おおむね静岡県の富士川・新潟県の糸魚川より東）が\qty{50}{Hz}、西日本が\qty{60}{Hz}です（電-10 ①）。1つの国の中に2つの周波数があるのは、世界でもめずらしいことです。

\phboxhead{明治の発電機が、今も東西を分けている}
電灯の会社が日本で発電を始めた明治の半ば（1890年代）、東京の会社はドイツ製の\qty{50}{Hz}の発電機を、大阪の会社はアメリカ製の\qty{60}{Hz}の発電機を入れました。そのころは、それぞれの町の中だけで電気を使っていたので、周波数が違っても困りませんでした。その後、東と西で、それぞれの周波数のまま電気の網が広がり、つながって大きくなりました。周波数をそろえる話は何度も出ましたが、発電機やモーターを全部取りかえる費用と、取りかえる間の停電を考えると、とても難しく、今も2つのままです。

\phboxhead{東と西を、直流でつなぐ}
周波数の違う2つの系統は、交流のまま電線でつなぐことができません。足並みのそろわない2つの大きな機械を、かたい軸で結ぶようなものだからです。そこで、\textbf{周波数変換所}で、片方の交流をいったん直流に変え（⑤の整流）、すぐにもう片方の周波数の交流にもどします（インバーター）。\bkidx{しゅうはすうへんかんしょ}{周波数変換所}直流には周波数がないので、東と西の発電機は、たがいの回る速さにしばられずに、電力だけをやりとりできるのです（図\ref{ph:電-16:fig-touzai}）。
\begin{bkfigure}
\begin{tikzpicture}[x=1mm, y=1mm]
  \sxqbox{hk}{(12,32)}{22}{北海道\\50\,Hz}{bkBlue}
  \sxqbox{ea}{(50,32)}{24}{東北・東京\\50\,Hz}{bkBlue}
  \sxqbox{we}{(96,32)}{30}{中部・北陸・関西\\中国・四国・九州\\60\,Hz}{bkOrange}
  \sxqbox{ok}{(108,13)}{20}{沖縄\\60\,Hz}{bkOrange}
  % 修4-6（G-8）：境の破線は「変換所」「210万 kW」の所で切り、ラベルの白地をやめた
  \draw[bkGray, line width=0.6pt, dashed] (73,7) -- (73,25.5) (73,38.5) -- (73,46);
  \draw[bkInk, line width=1.6pt] (hk) -- node[sxq lab, above=2pt, align=center] {直流} node[sxq lab, below=2pt, align=center] {90万 kW} (ea);
  \draw[bkInk, line width=1.6pt] (ea) -- node[sxq lab, above=2pt, align=center] {変換所} node[sxq lab, below=2pt, align=center] {210万 kW} (we);
  \node[sxq lab, anchor=north west] at (74,46) {60\,Hz};
  \node[sxq lab, anchor=north east] at (72,46) {50\,Hz};
  \node[sxq lab, anchor=west, align=left] at (0,13) {沖縄は、ほかとつながっていない};
\end{tikzpicture}
\bkfigcaption{日本の電力系統のつながり（形を見るための図。地図ではない）。北海道と本州は直流の送電線で、50\,Hz の東と 60\,Hz の西は周波数変換所（図では「変換所」）でつながる。数は、つないでやりとりできる電力の大きさ（2021年の時点）。}\label{ph:電-16:fig-touzai}
\end{bkfigure}
東と西の間の周波数変換所は、静岡県と長野県・岐阜県にあります。2011年の東日本大震災のあと、東日本で電気が足りなくなったとき、西から送れる電力が少なかったことから、増やす工事が進められました。2021年3月に、長野県と岐阜県を結ぶ設備（90万 kW）が運転を始め、東と西でやりとりできる電力は120万 kW から210万 kW に増えました。さらに300万 kW に増やす工事が進められています（完成の時期は、工事がおくれて見直されています）。210万 kW から300万 kW にすると、約43\,\% の増加です。

\begin{phmijika}[北海道と本州は、海の下の直流でつながる]
北海道と本州（青森県）の間も、直流の送電線でつながっています。もともと60万 kW の設備があり、2019年3月に、津軽海峡の下の鉄道のトンネルの中にケーブルを通した新しい設備（30万 kW）が加わって、90万 kW になりました。北海道も本州の東北も\qty{50}{Hz}ですが、海をこえる長いケーブルでは、⑥で見るように、直流のほうが都合がよいのです。
\end{phmijika}

%% ================================================================
\bksubsection{4}{大停電はどう広がるか：2018年の北海道}
\leavevmode\label{ph:電-16:sec4}②のしくみがあっても、いちどに大きく電源を失うと、守りきれないことがあります。2018年9月6日の\textbf{北海道胆振東部地震}では、日本で初めて、1つの系統の全体が停電する「ブラックアウト」が起きました。\bkidx{ぶらっくあうと}{ブラックアウト}経過は、電力広域的運営推進機関（国が認めた、全国の電力系統の運用を見る機関）の検証委員会の報告で、次のようにまとめられています（図\ref{ph:電-16:fig-hokkaido}）。
\begin{bkfigure}
\begin{tikzpicture}[x=1mm, y=1mm]
  \sxqbox{a}{(18,40)}{36}{3時7分　地震\\（M6.7、最大震度7）}{bkGray}
  \sxqbox{b}{(63,40)}{44}{火力発電所の2基が止まる\\（約116万 kW）。送電線の\\事故で水力も止まる}{bkRed}
  \sxqbox{c}{(112,40)}{36}{周波数が\\急に下がる}{bkRed}
  \sxqbox{d}{(112,18)}{46}{本州から直流で送る\\一部の地域を自動で\\切りはなす（約130万 kW）}{bkBlue}
  \sxqbox{e}{(63,18)}{44}{周波数は一時もちなおすが\\残った火力の1基も止まり\\ふたたび下がる}{bkRed}
  \sxqbox{f}{(18,18)}{36}{3時25分　ほぼ全域が\\停電（最大約295万戸）}{bkPurple}
  \draw[sxq arr] (a) -- (b); \draw[sxq arr] (b) -- (c); \draw[sxq arr] (c) -- (d); \draw[sxq arr] (d) -- (e); \draw[sxq arr] (e) -- (f);
\end{tikzpicture}
\bkfigcaption{2018年9月6日の北海道の全域停電までの流れ（電力広域的運営推進機関の検証委員会の資料から、要点だけをまとめた）。}\label{ph:電-16:fig-hokkaido}
\end{bkfigure}
地震の直後、震源の近くにあった大きな石炭火力発電所の2基が止まり、送電線の事故で水力発電所も止まりました。周波数は急に下がり、本州からの直流の送電を増やす装置と、周波数低下リレー（②）が働いて、合わせて約130万 kW の使う量を自動で切りはなしました。周波数は一時もちなおしましたが、残っていた同じ発電所の1基も止まり、ほかの発電所も次々に系統から切りはなされて、地震から18分後の3時25分に、北海道のほぼ全域が停電しました。

この経験から、周波数低下リレーで自動で切りはなす量を約35万 kW 増やすなどの対策がとられました。また、半年後の2019年3月に運転を始めた北海道と本州の新しい直流の設備（③）には、北海道の全体が停電していても、本州から電気を送り始めて停電を直すための機能（ブラックスタート）がついています。

\begin{phmada}[どの小さな故障が、大停電に育つのか？]
大停電は、ほとんどの場合、1つの故障だけでは起きません。1本の送電線が止まると、その分の電力がほかの線に流れこみ、それが重すぎて次の線が止まり……と、故障が次々に広がっていきます（連鎖）。北米の何十年分もの停電の記録を調べた研究では、停電の大きさの分布が、小さい停電ほど多く、大きい停電はとても少ないが「思ったより起きる」形（べき乗則）に近いことが示され、電力系統が、限界のふちで動くようにみずから近づいているのではないか、という考えが出されました。一方で、計算のモデルによっては、停電は「とても小さいか、とても大きいか」の2つに分かれる、という結果も出ていて、どちらがどんな系統で本当に起きているのかは、まだ決着していません。どの小さな故障が大停電に育つのかを、前もって見分けることも、まだ難しい問題です。
\end{phmada}

%% ================================================================
\bksubsection{5}{電気の形を変える：パワーエレクトロニクス}
\leavevmode\label{ph:電-16:sec5}発電所がつくるのは交流ですが、スマホ・パソコン・LED・電池は直流で動きます。電気自動車のモーターは、周波数を変えられる交流で回ります（電-12 ⑨）。このように、電気の形（直流か交流か、電圧、周波数）を、半導体のスイッチで変える技術を\textbf{パワーエレクトロニクス}といいます。\bkidx{ぱわーえれくとろにくす}{パワーエレクトロニクス}ここでも、スイッチは「オンかオフだけの理想の部品」、ダイオードは「一方の向きにだけ電流を通す弁」として扱います（半導体のしくみは物理シリーズの『原子』（原-9・原-10）で学びます）。
\begin{bkcenter}
\small
\begin{bktable}{colspec={Q[l,m,wd=20mm] Q[l,m,wd=36mm] X[l,m]}, width=\linewidth, rowsep=0.6pt, row{1}={bg=bkTheme!10, font=\sffamily\footnotesize}}
変える向き & 名前 & 使われている所 \\
交流 → 直流 & 整流器 & 充電器・パソコンの電源、直流送電の入口 \\
直流 → 交流 & インバーター & モーター（電-12 ⑨）、太陽光発電、直流送電の出口 \\
直流 → 直流 & DC-DC コンバーター & スマホの中、電気自動車の電池とモーターの間 \\
交流 → 交流 & 周波数変換 & 東西の周波数変換所（③） \\
\end{bktable}
\bktabcaption{電気の形を変える4つの向き。交流 → 交流は、ふつう、いったん直流にしてから交流にもどす。}\label{ph:電-16:tab-pe}
\end{bkcenter}

\phboxhead{整流：交流を一方の向きにそろえる}
ダイオードを4つ組み合わせると、交流の電圧がプラスのときも、マイナスのときも、出口では同じ向きに電流が流れるようにできます（\textbf{全波整流}）。\bkidx{せいりゅう}{整流}出口の電圧は、もとの正弦波のマイナスの半分を折り返した形 $\lvert v\rvert$ になります（図\ref{ph:電-16:fig-seiryu}）。ダイオード1つだけなら、マイナスの半分を捨てた形になります（半波整流）。
\begin{bkfigure}
\begin{tikzpicture}
\begin{axis}[bkplot, axis x line*=bottom, axis y line*=left, width=100mm, height=50mm, xmin=0, xmax=41, ymin=-160, ymax=170, xtick={0,10,20,30,40}, ytick={-141,0,90,141}, xlabel={時刻（ms）}, ylabel={電圧（V）}]
  \draw[bkGray, line width=0.5pt] (axis cs:0,0) -- (axis cs:41,0);
  \addplot[bkGray!80!black, line width=0.9pt, densely dotted, domain=0:40, samples=161] {141.4*sin(360*x/20)};
  \addplot[bkBlue!80!black, line width=1.5pt, domain=0:40, samples=161] {abs(141.4*sin(360*x/20))};
  \draw[bkOrange!85!black, line width=1.1pt, dashed] (axis cs:0,90) -- (axis cs:40,90);
  \node[sxq lab, anchor=west] at (axis cs:41.5,90) {平均};
\end{axis}
\end{tikzpicture}
\bkfigcaption{全波整流（\qty{50}{Hz}、実効値\qty{100}{V}）。もとの交流（点線）のマイナスの半分が折り返されて、実線のようにいつもプラスになる。破線が平均で、約\qty{90}{V}（\bkreidairef{2}）。}\label{ph:電-16:fig-seiryu}
\end{bkfigure}
整流しただけの電圧は、向きはそろっていますが、1秒に100回、0と\qty{141}{V}の間を上下しています。その平均の値は、5-12⑧の「面積を全体のはばでわる」で求められ、実効値\qty{100}{V}の交流では約\qty{90}{V}です（\bkreidairef{2}）。実効値（電-10 ②）の\qty{100}{V}とは違うことに注意しましょう。実効値は「熱を出すはたらき」でそろえた値、平均は「ならした高さ」です。

出口に大きなコンデンサーをつなぐと、電圧が高い間に充電され、低い間に放電して、でこぼこがならされます（平滑）。こうして、ほぼ一定の直流になります。パソコンの電源のスイッチを入れた瞬間に大きな電流が流れこむのは、この空のコンデンサーを充電するからでした（電-9 ⑪の突入電流）。

\phboxhead{DC-DC コンバーター：オンの割合で電圧を下げる}
直流の電圧を下げるには、スイッチを1秒に何万回もオン・オフし、オンの時間の割合を変えます。入口が\qty{400}{V}で、オンの時間が全体の4分の1なら、出口の電圧は、ならすと $400\times\tfrac14=\qty{100}{V}$ です（コイルとコンデンサーで、でこぼこをならす）。抵抗で電圧を下げるのと違い、スイッチはオンのときは電圧が0、オフのときは電流が0なので、理想ならスイッチの中で熱がまったく出ません。だから、パワーエレクトロニクスの装置は、効率がとても高いのです。電-12 ⑨のインバーターの PWM も、同じ「オンの割合で平均をつくる」考えです。

\Needspace{7\baselineskip}
\begin{phmijika}[ACアダプターが小さくなったわけ]
昔の充電器は、\qty{50}{Hz}の交流のまま変圧器で電圧を下げてから整流したので、大きく重い鉄心が要りました。今の充電器は、まず整流して直流にし、それを1秒に何万回以上もスイッチで切りかえた高い周波数の交流にしてから、小さな変圧器で電圧を下げます。電-12 ⑥の $V\approx4.44fN\Phi_{\max}$ から、同じ電圧なら、周波数が高いほど磁束が小さくてすみ、鉄心を小さくできるからです。
\end{phmijika}

\phboxhead{波の形のくずれと電磁ノイズ}
整流器やインバーターは、電流を細切れに取りこんだり、急に切りかえたりします。そのため、電流の波の形がくずれて、高調波（電-10 ⑪）をふくむようになり、変圧器をむだに熱くしたり、ほかの機械の動きを乱したりします。電気製品が出してよい高調波の量には決まりがあります。また、速い切りかえは電磁波の雑音（電磁ノイズ）も出すので（電-15 ⑫）、フィルターや遮へいで外にもらさないようにしています。

%% ================================================================
\bksubsection{6}{直流で送る}
\leavevmode\label{ph:電-16:sec6}発電所から家まで交流を使うのは、変圧器で電圧をかんたんに変えられるからでした（電-10 ①）。それでも、次のような場合には、交流をいったん直流にして送る\textbf{直流送電}が使われます。\bkidx{ちょくりゅうそうでん}{直流送電}
\begin{itemize}[leftmargin=1.2\zw, itemsep=2pt, topsep=2pt, label={\color{bkTheme}●}]
\item \textbf{周波数の違う系統をつなぐ}：東西の周波数変換所（③）。
\item \textbf{長い海底のケーブル}：ケーブルは、中心の電線と、外側を包む金属の層の間に、うすい絶縁物をはさんだ、長い長いコンデンサーです（電-3 ⑩の同軸の形）。交流で送ると、この静電容量を充電する電流（電-13 ⑩のフェランチ効果を起こす電流と同じもの）が、ずっと行ったり来たりし、それだけでケーブルの容量の多くを使ってしまいます。直流なら、はじめに1回充電するだけです。北海道と本州の連系（③）がこれです。
\item \textbf{とても長い距離}：送電線の数を減らせ、交流のような「足並みをそろえる」問題も起きません。
\end{itemize}
直流の送電線の入口と出口には、交流と直流を変える変換所が要ります（図\ref{ph:電-16:fig-hvdc}）。
\begin{bkfigure}
\begin{tikzpicture}[x=1mm, y=1mm]
  \sxqbox{a1}{(7,10)}{14}{交流の\\系統}{bkBlue}
  \sxqbox{c1}{(35,10)}{22}{変換所\\（整流）}{bkGray}
  \sxqbox{c2}{(87,10)}{28}{変換所\\（インバーター）}{bkGray}
  \sxqbox{a2}{(117,10)}{14}{交流の\\系統}{bkOrange}
  \draw[bkInk, line width=1.0pt] (a1) -- node[sxq lab, above=2pt] {三相} (c1);
  \draw[bkInk, line width=1.0pt] (c2) -- node[sxq lab, above=2pt] {三相} (a2);
  \draw[bkRed!80!black, line width=1.8pt] (c1.east) -- node[sxq lab, above=2pt] {直流 \qty{250}{kV} など} (c2.west);
  \draw[sxq arr] (54,0) -- (72,0);
  \node[sxq lab, anchor=north] at (63,-1.5) {電力の向き（逆にもできる）};
\end{tikzpicture}
\bkfigcaption{直流送電。入口で交流を直流に、出口で直流を交流に変える。スイッチの切りかえ方を変えれば、電力の向きも、送る量も、すばやく自由に変えられる。}\label{ph:電-16:fig-hvdc}
\end{bkfigure}
変換所のスイッチの切りかえ方で、送る電力の向きと大きさを、すばやく自由に決められます。2019年の北海道と本州の新しい設備は、日本で初めて\textbf{自励式}とよばれる変換器を使っています。それまでの国内の直流送電は、つながる交流の系統の電圧を使ってスイッチを切りかえる方式（他励式）でした。自励式は、自分でスイッチを入・切できる半導体を使うので、交流の側が停電していても、交流をつくり出して送れます。④のブラックスタートができるのは、このためです。

%% ================================================================
\bksubsection{7}{再生可能エネルギーとインバーター}
\leavevmode\label{ph:電-16:sec7}日本の発電電力量は1年に約1兆 kWh で、2024年度は約9900億 kWh でした。そのうち、水力をふくむ再生可能エネルギーは23.1\,\% で、2023年度の値では、その4割ほどが太陽光（全体の約10\,\%）です。

太陽光と風力の発電には、系統にとって大きな違いが2つあります。
\begin{enumerate}[label={(\arabic*)}, leftmargin=2.4\zw, itemsep=2pt, topsep=2pt]
\item \textbf{天気で出力が変わる}：雲が来れば太陽光は数分で大きく減り、夜は0です。そのぶん、ほかの発電所が出力を変えて、つくる量を合わせ続けなければなりません（②）。晴れて使う量の少ない休日の昼などには、太陽光がつくりすぎることもあり、日本でも2018年10月に九州で、離島以外では初めて、太陽光の出力を一時的におさえる\textbf{出力制御}が行われました。
\item \textbf{回る重い発電機がない}：太陽電池は直流をつくり、インバーター（⑤）で交流にして系統に送ります。風力も、多くはインバーターを通してつながります。インバーターの電源は、②の「回る運動エネルギー」を持たないので、こうした電源の割合が増えると、系統の慣性定数 $H$ が小さくなり、発電所が止まったときに周波数が速く下がるようになります（②の2倍にすると？）。
\end{enumerate}
\begin{bkfigure}
\begin{tikzpicture}
\begin{axis}[bkplot, axis x line*=bottom, axis y line*=left, width=100mm, height=50mm, xmin=0, xmax=24, ymin=0, ymax=1.15, xtick={0,6,12,18,24}, ytick=\empty, xlabel={時刻（時）}, ylabel={電力（形だけ）}]
  \addplot[bkGray!80!black, line width=0.8pt, domain=0:24, samples=97] {0.62+0.13*sin(15*(x-9))+0.18*exp(-((x-19)^2)/6)};
  \addplot[bkOrange!85!black, line width=1.1pt, dashed, domain=5:19, samples=57] {0.62*exp(-((x-12)^2)/9)};
  \addplot[bkBlue!80!black, line width=2.0pt, domain=0:24, samples=97] {0.62+0.13*sin(15*(x-9))+0.18*exp(-((x-19)^2)/6)-(x>5)*(x<19)*0.62*exp(-((x-12)^2)/9)};
\end{axis}
\end{tikzpicture}
\bkfigcaption{晴れた日の1日の電力（形を見るための図。数は書いていない）。細い実線が使う電力、破線が太陽光の発電。火力などのほかの発電所が受けもつ分（太い実線。細い実線から破線を引いたもの）は、昼に大きく減り、太陽光が減る夕方に急に増える。その間、ほかの発電所は、すばやく出力を変え続ける。}\label{ph:電-16:fig-ichinichi}
\end{bkfigure}
晴れた日には、ほかの発電所が受けもつ分は、図\ref{ph:電-16:fig-ichinichi}のように、昼に大きく減り、日がかたむく夕方に急に増えます。これに合わせて、火力・水力・揚水発電（⑧）が、すばやく出力を変えます。

②のしくみは、回る発電機がたくさんあることを前提にしてきました。そこで今、インバーターのスイッチの切りかえ方をくふうして、周波数が変わり始めたら、すぐに電池などから電力を出し入れし、回る発電機と同じようにふるまわせる技術（擬似慣性・グリッドフォーミング）の研究と導入が進んでいます。

\phboxhead{太陽光の広さの見積もり}
太陽電池の出力は、晴れた日の真昼に近い、1 m$^2$ あたり\qty{1000}{W}の日光が当たった条件で表示します。今の太陽電池のモジュールの変換効率は約20\,\% なので、1 m$^2$ で約\qty{200}{W}です。大きな火力発電所1基ぶんの\qty{1}{GW}（100万 kW）を、晴れた真昼につくるには、約\qty{5}{km^2}のパネルが要ります（\bkreidairef{3}）。ただし、太陽光は夜や雨の日には発電しないので、1年を通してならすと、出せる電気は定格の15〜18\,\% ほどです（熱-16 ⑤の設備利用率）。

%% ================================================================
\bksubsection{8}{電気をためる}
\leavevmode\label{ph:電-16:sec8}つくる量と使う量を合わせる、いちばん直接の方法は、電気をためておいて、余ったときにため、足りないときに出すことです。
\begin{itemize}[leftmargin=1.2\zw, itemsep=2pt, topsep=2pt, label={\color{bkTheme}●}]
\item \textbf{揚水発電}：電気の余る時間に、下の池の水をポンプで上の池へくみ上げ、足りない時間に水を落として発電します（熱-16 ⑤。水\qty{1}{m^3}を\qty{500}{m}上げて約\qty{1.4}{kWh}）。\bkidx{ようすいはつでん}{揚水発電}ポンプを回すモーターと発電機は、同じ同期機（電-12 ⑧）を、回る向きを変えて使い分けます。くみ上げてから発電するまでに、全体で約3割のエネルギーを失います（総合の効率は約70\,\%）。日本全体で、発電の出力の合計は約2700万 kW あります（2022年の時点）。回る重い発電機なので、②の慣性にもなります。
\item \textbf{蓄電池}：リチウムイオン電池などの大きな電池を、変電所や太陽光の発電所の近くに置きます。電池は直流なので、インバーター（⑤）を通して系統につなぎます。電力の出し入れを1秒もかからずに切りかえられるので、②の調速機よりも速く周波数を支えられます。
\item \textbf{水素}：余った電気で水を電気分解して水素をつくり、ためておいて、燃料電池（電-4 ⑨）や発電所で電気にもどします。長い期間ためられる方法として、研究と実験が進められています。
\end{itemize}
\noindent どの方法でも、ためて出すときにエネルギーの一部が熱になって失われます（エネルギーの保存）。季節をこえるほど長い期間、国全体を動かせるほどの電気を、安くためる方法は、まだ見つかっていません（熱-16 のまだわかっていないこと）。

%% ================================================================
\bksubsection{9}{系統を守る：雷・ショートと電気の安全}
\leavevmode\label{ph:電-16:sec9}電力系統では、どこかでときどき故障が起きます。原因として多いものの1つが、雷です。送電線に雷が落ちると、大きな電圧の波（\textbf{サージ}）が、光に近い速さで電線の上を走ります（電-9 のこの先の学び・電-13 ⑦の電線の上の波）。鉄塔のいちばん上の\textbf{架空地線}（電-11 ①）は、雷を先に受けとめて地面へ流し、電線に落ちにくくします。変電所の入口には、ふだんは電流を通さず、決まった電圧をこえたときだけ電流を地面へ逃がす\textbf{避雷器}を置いて、サージが変圧器などの高価な機器に入りこまないようにしています。

電線どうしが雷や木の枝などでつながる\textbf{短絡}（ショート）が起きると、ふだんの何倍もの大きな電流が流れます（電-13 ⑤）。電線や変電所には、電流と電圧をいつもはかっている\textbf{保護リレー}があり、故障を見つけると、すぐに\textbf{遮断器}に「切れ」と命令して、故障した区間だけを切りはなします。\bkidx{ほごりれー}{保護リレー}故障した所を、周りの健全な所に影響が広がる前に、すばやく正しく切りはなすことが、大停電（④）を防ぐ第一歩です。網の目（①）になっているので、1つの区間を切っても、ほかの道で電気を届けられます。

\begin{column}{電気の安全：送電線・電柱のそばで}
家の中の安全（感電・漏電・コードの熱）は電-4 ⑫で見ました。ここでは、家の外の電線のそばでの安全です。
\begin{itemize}[leftmargin=1.2\zw, itemsep=1pt, topsep=2pt, label={\color{bkTheme}●}]
\item \textbf{切れた電線・たれ下がった電線には、絶対に近づかない}。倒れた電柱や、電線に引っかかったたこ・風船にも触れない。見つけたら大人に知らせ、電力会社に連絡してもらう。
\item 送電線や電柱の近くでは、たこあげ・ドローン・ラジコンを飛ばさない。釣りざおを近づけない。電圧の高い線には、触れなくても、近づくだけで電気が飛んで感電することがある。
\item 災害で壊れたり水につかったりした太陽光パネルも、日光が当たっている間は発電しているので、近づかない。
\item 停電のあと、家を離れて避難するときは、分電盤のブレーカーを切っておく。電気がもどったとき、倒れた電気ストーブなどから火事になるのを防ぐ。
\end{itemize}
\end{column}

\begin{phsaki}{電力系統工学とパワーエレクトロニクス}
発電機の回り方と電力の流れを、何千もの式の連立方程式（電-5 ⑩）と微分方程式（電-9）で計算し、故障のあとも発電機が足並みをそろえて回り続けられるか（安定度）を調べるのが\textbf{電力系統工学}です。スイッチの切りかえ方で電気の形を自由に変える回路と、その制御を学ぶのが\textbf{パワーエレクトロニクス}です。スイッチになる半導体（ダイオード・トランジスタ）のしくみは物理シリーズの『原子』（原-9・原-10）で、原子力発電のしくみは原-15 で学びます。
\end{phsaki}


\begin{kotoba}\small
\bkword{電力系統}{発電所・送電線・変電所・配電線と、電気を使う人までをつないだ全体。網の目になっている。}
\bkword{慣性定数}{回る発電機の運動エネルギーが、定格の出力の何秒分かを表す数 $H$〔s〕。大きいほど周波数が変わりにくい。}
\bkword{調速機}{タービンの回る速さ（周波数）をはかり、遅くなると蒸気や水を多く入れて出力を増やす装置。}
\bkword{周波数変換所}{交流をいったん直流にして、違う周波数の交流にもどし、東と西の系統をつなぐ所。}
\bkword{整流}{交流を、一方の向きにだけ流れる電気（直流）に変えること。}
\bkword{パワーエレクトロニクス}{半導体のスイッチで、電気の形（直流・交流、電圧、周波数）を変える技術。}
\bkword{直流送電}{交流を直流に変えて送り、受ける側で交流にもどす送電。周波数の違う系統・海底のケーブルに使う。}
\bkword{保護リレー}{電線の電流や電圧を見張り、故障を見つけると遮断器に切りはなす命令を出す装置。}
\end{kotoba}
\bkkigocards{\bkkigocard{$H$}{エイチ}{慣性定数〔s〕。回る発電機の運動エネルギーが、定格の出力の何秒分か。第2巻のエンタルピー、電-7 の磁場の強さ $H$ とは別}{inertia constant}{電-16}}{\bkkigocard{$P_{\mathrm n}$}{ピー・エヌ}{系統につながっている発電機の、定格の出力の合計〔W〕。電-16 では、つくる電力 $P_{\text{発}}$・使う電力 $P_{\text{使}}$ と分けて書く}{rated power}{電-16}}

%% ================================================================
\Needspace{12\baselineskip}\bksection{reidai}{見本の計算}

\begin{phreidai}[発電所が止まった直後の周波数]\label{ph:電-16:ex1}
\qty{50}{Hz}の系統で、つながっている発電機全体の定格の出力の1\,\% にあたる発電所が、急に止まりました。系統の慣性定数を $H=\qty{5}{s}$ とします。止まった直後、周波数は1秒に何 Hz の速さで下がり始めますか。また、そのまま何もしなければ、何秒で\qty{0.2}{Hz}下がりますか。
\phbamen
止まった発電所の分だけ、つくる電力が使う電力より少なくなり、足りない分を回る発電機の運動エネルギーが出します。
\phmushi
調速機などが働く前の、はじめの短い時間だけを考えます。使う電力が周波数で変わることを無視します。系統の発電機はすべてそろって回るとします。
\phzu
基準の周波数 $f_0=\qty{50}{Hz}$、慣性定数 $H=\qty{5}{s}$。\\
足りない電力の割合 $(P_{\text{使}}-P_{\text{発}})/P_{\mathrm n}=0.01$。
\phriyu
回る運動エネルギーの減る速さが足りない電力に等しい（②の式）。$K=K_0(f/f_0)^2$、$K_0=HP_{\mathrm n}$ から導いた②の流れ図の右の式を使います。
\phtateru
\[\frac{\dd f}{\dd t}=\frac{f_0}{2H}\cdot\frac{P_{\text{発}}-P_{\text{使}}}{P_{\mathrm n}}\]
\phtoku
\[\frac{\dd f}{\dd t}=\frac{50}{2\times5}\times(-0.01)=-\qty{0.05}{Hz/s},\qquad \text{時間}=\frac{0.2}{0.05}=\qty{4}{s}\]
\phtashikame
単位：Hz ÷ s ×（単位なし）＝ Hz/s。\\
大きさ：4秒で目標の幅（上下\qty{0.2}{Hz}）の外に出るので、調速機が数秒のうちに働く必要があります（②）。\\
2倍にすると：$H$ が半分（\qty{2.5}{s}）なら\qty{0.1}{Hz/s}で、2倍速く下がります。止まる発電所が2\,\% でも同じです。
\kotae 1秒に約\qty{0.05}{Hz}の速さで下がり始め、何もしなければ約4秒で\qty{0.2}{Hz}下がる
\end{phreidai}

\begin{phreidai}[整流した電圧の平均]\label{ph:電-16:ex2}
実効値\qty{100}{V}の交流（正弦波）を全波整流しました。出口の電圧の平均の値を求めましょう。半波整流ならいくらですか。
\phbamen
充電器の入口で、交流を一方の向きにそろえます（⑤）。
\phmushi
ダイオードの中で下がる電圧を無視し、ダイオードは理想の弁とします。平滑のコンデンサーはつながないとします。
\phzu
山の電圧 $V_0=\sqrt2\times100\approx\qty{141.4}{V}$。全波整流の出口の電圧 $\lvert V_0\sin\omega t\rvert$。同じ形が半周期ごとに繰り返すので、半周期（$0\le\omega t\le\pi$）で平均すればよい。
\phriyu
平均の値は、グラフの下の面積をはばでわったもの（5-12⑧）。$\omega t=\theta$ とおいて、$\theta$ で平均します。
\phtateru
\[\text{平均}=\frac{1}{\pi}\int_0^{\pi}V_0\sin\theta\,\dd\theta\]
\phtoku
\[\int_0^{\pi}\sin\theta\,\dd\theta=\Bigl[-\cos\theta\Bigr]_0^{\pi}=2,\qquad \text{平均}=\frac{2V_0}{\pi}=\frac{2\times141.4}{3.14}\approx\qty{90}{V}\]
半波整流では、残りの半周期は0なので、1周期で平均すると半分の約\qty{45}{V}です。
\phtashikame
単位：V × rad ÷ rad ＝ V。\\
大きさ：山の\qty{141}{V}より小さく、0より大きい。実効値\qty{100}{V}とも違う（平均は「ならした高さ」、実効値は「熱のはたらき」でそろえた値）。\\
2倍にすると：\qty{200}{V}なら、平均も2倍の約\qty{180}{V}。
\kotae 全波整流で約\qty{90}{V}、半波整流で約\qty{45}{V}
\end{phreidai}

\begin{phreidai}[太陽光で 1 GW をつくる広さ]\label{ph:電-16:ex3}
変換効率20\,\% の太陽電池のパネルに、1 m$^2$ あたり\qty{1000}{W}の日光が当たるとします。\qty{1}{GW}（100万 kW）の電力をつくるには、パネルの面積はいくら要りますか。
\phbamen
大きな火力発電所1基ぶんの電力を、晴れた真昼の太陽光でつくる場面です（⑦）。
\phmushi
パネルの温度で効率が下がること、インバーターの損失、パネルどうしのすき間、影を無視します。日光はパネルに垂直に当たるとします。
\phzu
日光の強さ $G=\qty{1000}{W/m^2}$、効率 $\eta=0.20$、つくる電力 $P=\qty{1.0e9}{W}$、面積 $A$。
\phriyu
1 m$^2$ でつくれる電力は $\eta G$。全体の電力は、それに面積をかけたものです。
\phtateru
\[P=\eta GA\qquad\text{だから}\qquad A=\frac{P}{\eta G}\]
\phtoku
\[A=\frac{1.0\times10^9}{0.20\times1000}=\qty{5.0e6}{m^2}=\qty{5.0}{km^2}\]
\phtashikame
単位：W ÷（W/m$^2$）＝ m$^2$。\\
大きさ：一辺が約\qty{2.2}{km}の正方形です。すき間を入れると、もっと広くなります。\\
変えると：効率が10\,\% なら2倍の\qty{10}{km^2}。夜や曇りの日があるので、1年でならした電気の量で火力1基と並ぶには、さらに何倍もの広さが要ります（熱-16 ⑤）。
\kotae 約\qty{5}{km^2}（一辺 約\qty{2.2}{km}の正方形）
\end{phreidai}

\begin{phreidai}[直流で北海道へ送る]\label{ph:電-16:ex4}
2019年に運転を始めた、北海道と本州を結ぶ新しい直流の送電の設備は、電圧\qty{250}{kV}・電流\qty{1200}{A}で送ります。送る電力はいくらですか。送電線（片道）の抵抗を\qty{2}{\ohm}とすると、線の損失は送る電力の何\,\% ですか。同じ電力を半分の\qty{125}{kV}で送ると、損失は何\,\% になりますか。
\phbamen
直流で、海をこえて電力を送る場面です（③⑥）。
\phmushi
変換所の中の損失を無視します。帰りの電流の道（大地など）の抵抗は、線の抵抗にふくめて考えます。抵抗\qty{2}{\ohm}は仮の値です。
\phzu
電圧 $V=\qty{2.5e5}{V}$、電流 $I=\qty{1200}{A}$、線の抵抗 $R=\qty{2}{\ohm}$、送る電力 $P$、損失 $P_{\text{損}}$。
\phriyu
直流では、電力は $P=VI$（電-4 ⑤）、線の損失は $RI^2$ です。交流のように力率（電-11 ③）を考えなくてよいのも、直流の便利な所です。
\phtateru
\[P=VI,\qquad P_{\text{損}}=RI^2,\qquad \text{割合}=\frac{P_{\text{損}}}{P}\]
\phtoku
\[P=2.5\times10^5\times1200=\qty{3.0e8}{W}=\qty{300}{MW}\]
\[P_{\text{損}}=2\times1200^2\approx\qty{2.9e6}{W}\]
割合は $2.9\div300\approx0.010$、約1\,\% です。\qty{125}{kV}で送ると、電流は2倍の\qty{2400}{A}、損失は $2\times2400^2\approx\qty{1.2e7}{W}$ で、約4\,\% です。
\phtashikame
単位：V × A ＝ W、Ω × A$^2$ ＝ W。\\
大きさ：\qty{300}{MW}（30万 kW）は、③の「新しい設備 30万 kW」と合います。\\
2倍にすると：電圧を半分にすると、電流が2倍で、損失は4倍。電-11 ⑥の「損失 $R(P/V)^2$」と同じです。
\kotae 送る電力 \qty{300}{MW}。損失は約1\,\%、\qty{125}{kV}なら約4\,\%
\end{phreidai}

%% ================================================================
\bksection{tsumazuki}{まちがえやすいところをチェック}
\begin{tsumazuki}{電気はためて使っている？}
\bkng{発電所は、つくった電気を大きなタンクにためておき、使う分だけ送っている。}
\bkok{電気はふつう、つくったその瞬間に使われる。だから、つくる量を使う量に合わせ続ける。}
ためる設備（⑧）もありますが、全体から見るとわずかです。合っているかどうかを、周波数で見張っています（②）。
\end{tsumazuki}
\begin{tsumazuki}{東の電気を西へ}
\bkng{東日本で余った電気は、送電線をつなげば、そのまま西日本へ送れる。}
\bkok{周波数が違うので、周波数変換所で直流にしてから、60\,Hz の交流にもどして送る。送れる量にも上限がある。}
2021年の時点で、東と西でやりとりできるのは210万 kW までです（③）。
\end{tsumazuki}
\begin{tsumazuki}{整流した電圧}
\bkng{\qty{100}{V}の交流を全波整流すると、平均\qty{100}{V}の直流になる。}
\bkok{平均は約\qty{90}{V}。\qty{100}{V}は実効値で、平均とは違う値。}
平均は「ならした高さ」、実効値は「熱を出すはたらき」でそろえた値です（\bkreidairef{2}）。コンデンサーでならすと、山の\qty{141}{V}に近づきます。
\end{tsumazuki}

%% ================================================================
\Needspace{12\baselineskip}\bksection{phgear}{工学ではこう使う}
\begin{phkougaku}[次の日の電気の量を予想する]
中央給電指令所では、前の日に、次の日の30分ごとの使う量を予想します。夏の暑い日の午後はエアコンで、冬の寒い朝は暖房で、使う量が増えます。休日は工場が休むので、平日より少なくなります。太陽光が多くなった今は、天気予報から、太陽光がどれだけ発電するかも予想します（図\ref{ph:電-16:fig-ichinichi}）。

予想した量に、急な故障や予想のはずれに備えた余裕（予備力）を足して、どの発電所を何時から動かすかを決めます。余裕が小さくなりそうなときは、ほかの地域から送ってもらったり（③）、使う人に節電をお願いしたりします。「でんき予報」の使用率は、この余裕の目安を、だれでも見られるようにしたものです。当日、予想がはずれた分は、②の調速機と周波数をもとにもどす制御が、その場で合わせます。
\begin{phketa}[電気を届ける数]
\phketaline{周波数の目標}{上下\qty{0.2}{Hz}（北海道は上下\qty{0.3}{Hz}）}
\phketaline{東西の連系／北海道と本州}{210万 kW／90万 kW（直流）}
\phketaline{発電電力量（2024年度）}{約9900億 kWh（再エネ 23.1\,\%）}
\phketaline{揚水発電（2022年の時点）}{合計 約2700万 kW、総合の効率 約70\,\%}
\phketaline{慣性定数・太陽電池のモジュール}{数秒・効率 約20\,\%}
\phboxhead{答えが桁で合っているか}
周波数のずれは、ふだんは\qty{0.1}{Hz}ほど。1秒に何 Hz も下がる答えが出たら、\% と小数の取り違えや、$f_0/(2H)$ の2を確かめます。kW と万 kW の取り違えにも気をつけます。
\end{phketa}
\end{phkougaku}
\phstation{電力系統工学,パワーエレクトロニクス,電気設備,高電圧工学}{需給バランス、周波数制御（ガバナフリー・LFC・EDC）、慣性定数、周波数低下リレー、周波数変換所、直流送電（HVDC）、整流器・インバーター・DC-DC コンバーター、保護リレー、避雷器}

\begin{furikaeri}
\begin{checklist}
\item 発電所から家までの電気の道すじと、東と西で周波数が違うわけ・東西をつなぐしくみを説明できる。\bkfurref{ph:電-16:sec1}
\item 整流・インバーター・直流送電を、スイッチの切りかえと「平均の値」で説明し、整流した電圧の平均を計算できる。\bkfurref{ph:電-16:ex2}
\item つくる電気と使う電気がずれると周波数が変わるわけを、回る発電機の運動エネルギーで説明し、周波数の下がる速さを見積もれる。\bkfurref{ph:電-16:sec2}
\item 調速機の働きを入れた周波数の変わり方を、プログラムで計算して読める。\bkfurref{ph:電-16:prog}
\end{checklist}
\end{furikaeri}
\bkfukidashi{mina}{電気の網の中では、何百台もの発電機がそろって回り、周波数がつり合いを知らせている。場の法則から始めたこの巻の道具が、コンセントの向こうで全部はたらいているんだね。}
